\section{Burst Splitter}
\label{sec:splitter}
\subsection{Purpose and bounds}

The Burst Splitter converts one accepted AXI transaction into one to four MC
packet segments.  It is shared by the read and write request paths.  A row
boundary is the only splitting boundary: AXI beats are carried by a segment's
\texttt{beat\_cnt}; they are not independently converted into MC packets.

The configured \texttt{row\_size} is a run-time parameter and is at least
1~KB.  AXI4 prohibits a burst from crossing a 4~KB boundary, so a transaction
spans at most 4~KB.  It can therefore cross at most three 1~KB rows and has
at most four segments:
\[
 \texttt{MAX\_SEGMENTS}=1+3=4.
\]

\subsection{Interfaces and state}

At transaction acceptance, the selected read or write ROB allocates
\texttt{ROB\_INDEX[7:0]}.  The splitter accepts that index with the AXI
metadata and retains it for the whole transaction.  Stage~1 determines
\texttt{num\_segments} and writes the Segment Table.  It presents
\texttt{splitter\_pkt\_count[2:0]=num\_segments} to the corresponding ROB;
this is the sole source of \texttt{exp\_pkts}.  It is not derived from
\texttt{AWLEN+1} or \texttt{ARLEN+1}.

\begin{center}
\begin{tabular}{lll}
\toprule
\textbf{State} & \textbf{Width/depth} & \textbf{Meaning}\\
\midrule
Segment Table & 4 entries & \texttt{\{addr[31:0], beat\_cnt[7:0]\}} per segment\\
\texttt{ROB\_INDEX} & 8 bits & Transaction index retained while segments issue\\
\texttt{num\_segments} & 3 bits & Number of valid Segment Table entries, 1--4\\
\texttt{seg\_idx} & 2 bits & Segment currently being emitted\\
\bottomrule
\end{tabular}
\end{center}

\subsection{Stage 1 --- segment construction}

Stage~1 receives address, burst length, size, burst type, AXID, packet type,
and \texttt{ROB\_INDEX}.  Using the run-time \texttt{row\_size}, it walks
the addressed burst span and partitions it only at row boundaries.  It writes
one Segment Table entry for each contiguous portion of the transaction.  The
resulting count is in the inclusive range 1--4.  The AXI 4~KB constraint and
minimum row size make a fifth entry unreachable.

For a read allocation, Stage~1 also provides the initialized count to the
read ROB before its \texttt{ROB\_INDEX} becomes visible at the relevant
per-AXID Slot FIFO.  The write path follows the identical ordering.

\subsection{Stage 2 --- sequential emission}

Stage~2 emits Segment Table entries in increasing \texttt{seg\_idx} order.
For each accepted packet it forms
\[
 \texttt{GLOBAL\_TAG[9:0]=\{ROB\_INDEX[7:0],PACKET\_NUMBER[1:0]\}},
 \qquad \texttt{PACKET\_NUMBER=seg\_idx}.
\]
Packet Formation receives \texttt{addr}, \texttt{beat\_cnt}, packet type,
and \texttt{GLOBAL\_TAG}; it serializes one MC request packet per Segment
Table entry.  For reads, the packet-issue path also obtains a pooled Read Data
Buffer pointer and records it in \texttt{dbuf\_ptr[seg\_idx]} before that
packet can complete.

Stage~2 advances \texttt{seg\_idx} only on downstream packet acceptance.  It
stalls Stage~1 whenever it is emitting a multi-segment transaction, thereby
preserving the Segment Table and its associated \texttt{ROB\_INDEX}.  After
the final accepted segment it returns to idle and permits the next Stage~1
transaction.

\subsection{WRAP addressing within the Segment Table model}

The Segment Table and WRAP address generation solve different problems and
operate together for a WRAP burst that crosses a row boundary:

\begin{itemize}[noitemsep]
  \item Stage~1's row-boundary partitioning applies identically to INCR and
        WRAP bursts: it determines \emph{how many} row-bounded segments the
        transaction requires and each segment's \texttt{addr}/\texttt{beat\_cnt}
        starting point, independent of burst type.
  \item WRAP address generation determines the \emph{per-beat} address
        sequence \emph{within} a segment.  Because a WRAP burst's beat
        addresses are not simply contiguous beyond its wrap boundary, the
        beat-by-beat address a WRAP segment presents to MC Core is produced
        by WRAP address generation, not by incrementing \texttt{addr} by the
        transfer size each beat.
  \item Both mechanisms are required together for a WRAP burst that crosses
        one or more row boundaries: the Segment Table determines segment
        count and each segment's starting point, while WRAP address
        generation determines the beat addresses used within each segment.
        Neither mechanism substitutes for the other.
\end{itemize}

This preserves the existing WRAP Beat Address Table semantics --- a
sequential, per-beat address sequence consumed in beat order --- scoped now
to one Segment Table entry at a time rather than to the whole burst.

\subsection{Backpressure and error packets}

\begin{itemize}[noitemsep]
  \item A full selected ROB Slot FIFO, an empty selected ROB Free List, or
        unavailable read Data Buffer capacity prevents acceptance upstream as
        appropriate to that path.
  \item Packet-FIFO backpressure holds the current Segment Table entry and
        its \texttt{seg\_idx}; it does not alter \texttt{num\_segments}.
  \item Error read and error write transactions use the same segmentation and
        \texttt{GLOBAL\_TAG} mechanism.  Their packet type is carried in
        every generated segment.
\end{itemize}

\subsection{Timing notes}

\begin{itemize}[noitemsep]
  \item \texttt{splitter\_pkt\_count} is the number of packets generated by
        this one transaction, not a sum across a burst.
  \item \texttt{seg\_idx} is assigned at packet emission, not allocation.
  \item Segment Table state is indexed only by \texttt{seg\_idx}; packet
        counting is independent of AXI beat count.
\end{itemize}
